\chapter{Manutenção preventiva e corretiva}
\label{cap:manutencao}

\begin{objetivos}
Ao final deste capítulo você deve ser capaz de:
\begin{itemize}
  \item reconhecer os modos de falha típicos de um sistema com CLP;
  \item montar uma rotina preventiva com período e evidência para cada item;
  \item diagnosticar por sintoma, seguindo um roteiro de verificação;
  \item executar uma intervenção em campo com bloqueio e registro;
  \item identificar quando a manutenção preditiva é viável e quando é apenas dado sem uso.
\end{itemize}
\end{objetivos}

\section{Três formas de manter}

\begin{itemize}
  \item \textbf{corretiva}: age depois da falha. É inevitável em algum grau, mas cara: a parada não é
        planejada, o sobressalente pode não existir e a produção está perdida;
  \item \textbf{preventiva}: age em intervalo programado, substituindo ou verificando antes de a falha
        acontecer. Custa mesmo quando o componente estaria bom --- mas a parada é planejada;
  \item \textbf{preditiva}: age quando a condição medida indica degradação. Reduz intervenção
        desnecessária, mas exige medição, histórico e critério --- sem os três, vira preventiva com
        nome novo.
\end{itemize}

Em automação, o ganho maior não está em escolher uma das três, e sim em \emph{deslocar o esforço para
o lado barato}: registro de falha bem feito (que permite o diagnóstico rápido), sobressalente
disponível e intervenção com procedimento escrito.

\index{manutenção!preventiva}\index{manutenção!corretiva}\index{manutenção!preditiva}

\section{Modos de falha típicos}

\begin{table}[H]
  \centering
  \caption{Modos de falha usuais em sistemas com CLP, sinal observável e prevenção.}
  \label{tab:modos-falha}
  \begin{tabular}{p{3.2cm} p{4.6cm} p{3.2cm} p{3.3cm}}
    \toprule
    Componente & Modo de falha & Sinal observável & Prevenção \\
    \midrule
    Bateria de retenção
      & Descarga: as variáveis retentivas e o relógio de tempo real são perdidos quando a energia cai
      & Perda de contadores e de horário após parada; aviso de bateria fraca no diagnóstico do
        controlador
      & Substituição em intervalo definido, com registro da data, e conferência do aviso de
        diagnóstico (Capítulo~\ref{cap:memoria}) \\
    \addlinespace
    Cartão de entrada ou saída
      & Canal queimado por surto, curto ou sobrecarga
      & Um ponto específico não lê ou não atua, enquanto os demais do mesmo cartão funcionam
      & Proteção de circuito, separação de cargas indutivas, conferência dos sinais de campo e
        registro de falha por ponto \\
    \addlinespace
    Fonte de alimentação
      & Envelhecimento de capacitores, sobretensão, sobrecarga
      & Reinícios aleatórios, falha intermitente de módulos, reinício em partida de motor
      & Medição de tensão sob carga, dimensionamento correto, reserva de capacidade e ventilação do
        painel \\
    \addlinespace
    Térmica e poeira
      & Acúmulo de pó, filtro saturado, ventilação obstruída
      & Falha que aparece no verão ou depois de horas de operação e desaparece depois de esfriar
      & Limpeza e troca de filtro em rotina; medição de temperatura no interior do painel \\
    \addlinespace
    Programa e memória
      & Corrupção de programa, bateria fraca, edição indevida em campo
      & Comportamento diferente do documentado, sem falha de hardware
      & Backup versionado, restrição de acesso e controle de quem altera o quê
        (Capítulos~\ref{cap:documentacao} e~\ref{cap:seguranca}) \\
    \addlinespace
    Rede de comunicação
      & Cabo, conector, terminador, endereço duplicado, interferência
      & Comunicação cai e volta; erro em um nó específico
      & Conferência de cabeamento e de aterramento, documentação do endereçamento e teste de
        continuidade \\
    \addlinespace
    Campo
      & Sensor descalibrado, atuador travado, fim de curso desregulado
      & Leitura plausível e errada --- o pior tipo de falha
      & Conferência periódica com instrumento padrão e teste de atuação
        (Capítulo~\ref{cap:simulacao}) \\
    \bottomrule
  \end{tabular}
\end{table}

\section{Rotina preventiva}

A rotina precisa ser executável por quem está no turno e verificável por quem responde pela planta.
Cada item tem período, responsável e evidência.

\begin{table}[H]
  \centering
  \caption{Rotina preventiva mínima e evidência de execução.}
  \label{tab:rotina-preventiva}
  \begin{tabular}{p{4.4cm} p{3.0cm} p{6.9cm}}
    \toprule
    Item & Período & Evidência \\
    \midrule
    Backup do projeto do controlador & A cada alteração e em intervalo fixo
      & Arquivo versionado, com identificação de versão do programa e data; cópia guardada fora do
        equipamento \\
    Leitura do diagnóstico do controlador & Mensal
      & Registro dos avisos ativos, incluindo bateria, temperatura e estatística de erros de
        comunicação \\
    Inspeção visual e limpeza do painel & Semestral ou conforme ambiente
      & Lista assinada: pó, ventilação, filtros, conexões, sinais de aquecimento \\
    Medição de temperatura interna & Semestral, nos meses quentes
      & Valor registrado, comparado ao histórico \\
    Teste de bateria de retenção & Conforme o fabricante e o histórico
      & Data de substituição anotada no painel e no sistema de manutenção \\
    Conferência de pontos de E/S & Anual, ou após intervenção
      & Lista de pontos testados, com resultado (Capítulo~\ref{cap:es}) \\
    Teste dos intertravamentos e da parada de emergência & Anual e após qualquer alteração de lógica
      & Roteiro de teste assinado, com data e responsável (Capítulo~\ref{cap:simulacao}) \\
    Restauração de backup em bancada & Anual
      & Registro do teste de restauração: sem isso, o backup é apenas uma esperança \\
    Revisão da documentação e da lista de tags & Anual
      & Documentação atualizada e publicada na versão vigente \\
    \bottomrule
  \end{tabular}
\end{table}

\begin{atencao}
Teste de intertravamento e de parada de emergência só pode ser feito com a máquina em condição
segura, com a área liberada e com autorização de quem responde pela produção. Testar proteção com
pessoa em zona de risco é a receita do acidente que a proteção existia para evitar
(Capítulo~\ref{cap:simulacao}).
\end{atencao}

\section{Diagnóstico por sintoma}

O roteiro da Figura~\ref{fig:fluxo-diagnostico} parte do sintoma e desce pela cadeia, sempre do mais
provável e mais barato de verificar para o menos provável. A ordem importa: conferir o programa antes
de conferir a fonte é o erro mais comum e o mais caro em tempo.

\begin{figure}[H]
  \centering
  \includegraphics[width=0.85\linewidth]{24-diagnostico.pdf}
  \caption{Roteiro de verificação a partir do sintoma observado.}
  \label{fig:fluxo-diagnostico}
\end{figure}

\begin{table}[H]
  \centering
  \caption{Diagnóstico: sintoma, causas prováveis e verificação.}
  \label{tab:diagnostico}
  \begin{tabular}{p{3.3cm} p{5.1cm} p{5.9cm}}
    \toprule
    Sintoma & Causas prováveis & Como verificar \\
    \midrule
    Saída não atua
      & Ponto forçado no ambiente de programação; canal do cartão queimado; lógica não satisfeita;
        circuito de campo aberto; fusível ou relé
      & Conferir a força ativa (\emph{force}); observar o LED do canal; verificar a condição que
        deveria acionar a saída; medir tensão no borne com o ponto comandado \\
    \addlinespace
    Entrada não lê
      & Sinal de campo ausente; tensão de alimentação incorreta; filtro de entrada longo demais para o
        evento; endereço errado na lista de tags; borneira solta
      & Medir o sinal no borne; conferir o tipo de entrada (sourcing ou sinking); comparar o tempo do
        evento com o filtro configurado; conferir a alocação do ponto \\
    \addlinespace
    CLP em \emph{fault}
      & Falha de fonte; módulo com erro; programa incompatível com o \emph{firmware}; erro de
        configuração de rack
      & Ler o diagnóstico do controlador; conferir tensões; testar módulo por módulo; verificar a
        versão do ambiente de programação \\
    \addlinespace
    Comportamento diferente do documentado
      & Programa alterado em campo; variável retentiva perdida; temporizador com base de tempo
        diferente da prevista
      & Comparar o programa carregado com o backup versionado; conferir bateria e retenção; medir o
        tempo real de operação com cronômetro \\
    \addlinespace
    Comunicação cai e volta
      & Cabo, conector, terminador, endereço duplicado, aterramento, interferência de acionamento
      & Conferir endereçamento e terminadores; revisar aterramento e roteamento dos cabos junto aos
        cabos de potência \\
    \addlinespace
    Leitura plausível, mas errada
      & Sensor descalibrado ou com sinal de campo degradado; escala ou unidade trocadas
      & Comparar com instrumento padrão; conferir escala e unidade na lista de tags
        (Capítulo~\ref{cap:documentacao}) \\
    \bottomrule
  \end{tabular}
\end{table}

\section{Procedimento de intervenção}

Este procedimento vale para qualquer intervenção em máquina em produção. Ele não é uma sugestão de
boa prática: é a sequência mínima que separa a manutenção do acidente.

\begin{enumerate}
  \item \textbf{comunicar} quem opera e quem responde pela área antes de começar;
  \item \textbf{bloquear e etiquetar a energia}: desligar, bloquear o dispositivo de seccionamento e
        identificar o responsável. Energia residual --- pressão, mola, gravidade, capacitor, inércia
        --- precisa ser dissipada ou contida. \emph{Nenhum trabalho de ligação, medição ou troca de
        componente é feito em equipamento energizado};
  \item \textbf{conferir a ausência de tensão} com instrumento adequado, no ponto onde se vai
        trabalhar, antes de tocar no circuito;
  \item \textbf{fazer backup antes de alterar} qualquer coisa no programa, e anotar a versão;
  \item \textbf{se a alteração for no programa}, usar o modo de operação mais seguro possível e
        manter a lógica de proteção ativa. Forçar ponto para ``testar'' com a máquina em produção é
        proibido: a força desabilita a proteção;
  \item \textbf{testar após a intervenção} pelos casos críticos do roteiro
        (Capítulo~\ref{cap:simulacao}), antes de liberar a produção;
  \item \textbf{registrar}: o que foi feito, por quê, quando, por quem, qual a versão do programa
        resultante e qual sobressalente foi consumido. Registro é o que transforma a próxima
        ocorrência em diagnóstico rápido, em vez de investigação do zero;
  \item \textbf{liberar} com a área segura e as proteções restabelecidas, comunicando quem
        bloqueou.
\end{enumerate}

\begin{atencao}
Dois hábitos causam acidente com controlador: forçar saída para testar com a máquina em operação e
editar o programa \emph{on-line} sem conferir o efeito imediato na máquina. A edição em modo de
execução altera o comportamento no ciclo seguinte --- o que significa que o efeito é instantâneo. Se
a edição for inevitável, faça-a com o equipamento em condição segura e com o intertravamento
intacto.
\end{atencao}

\section{Manutenção preditiva como evolução}

Predição útil não exige sensor exótico. Três fontes de dado simples, já presentes ou de fácil
obtenção, cobrem boa parte dos casos:

\begin{itemize}
  \item \textbf{contagem de ciclos e de manobras} por atuador, comparada à vida declarada pelo
        fabricante. É o dado mais confiável e o mais barato de manter
        (Capítulo~\ref{cap:dados});
  \item \textbf{temperatura}, no painel e no motor, com tendência comparada ao histórico. Degradação
        aparece como desvio lento, não como limite atingido;
  \item \textbf{corrente e tempo de etapa}: corrente de motor que sobe lentamente indica desgaste ou
        obstrução; tempo de etapa que cresce indica perda de desempenho
        (Capítulo~\ref{cap:eficiencia}).
\end{itemize}

O critério de viabilidade é econômico e operacional: só compensa predição quando o custo da falha
não planejada é maior que o custo de medir, guardar e analisar o dado --- e quando existe alguém
encarregado de olhar o resultado. Indicador de degradação publicado e não acompanhado é a pior forma
de preventiva: parece controle e não é.

\begin{pratica}
Escolha um equipamento automatizado da sua planta ou bancada e monte o plano de manutenção dele:
(a) liste os cinco modos de falha mais prováveis, no formato da Tabela~\ref{tab:modos-falha};
(b) preencha a Tabela~\ref{tab:rotina-preventiva} com período e responsável para cada item aplicável;
(c) redija o procedimento de intervenção em uma página, na ordem dos oito passos;
(d) descreva dois sintomas e o roteiro de verificação de cada um, usando o formato da
Tabela~\ref{tab:diagnostico}.
\end{pratica}

\begin{exercicios}
\begin{enumerate}[label=\alph*)]
  \item Explique a diferença entre manutenção preventiva e preditiva, e diga o que a preditiva exige
        além da medição.
  \item Por que a perda de variáveis retentivas após uma queda de energia é um modo de falha
        previsível?
  \item Uma saída não atua. Descreva o roteiro de verificação na ordem correta, com base na
        Tabela~\ref{tab:diagnostico}.
  \item Uma entrada não lê um pulso curto, mas lê um sinal contínuo. Qual é a causa mais provável e
        como confirmar?
  \item Por que a falha mais perigosa é a leitura plausível e errada? Dê um exemplo.
  \item Liste, em ordem, os oito passos do procedimento de intervenção e explique a função de cada um.
  \item Por que não se deve forçar ponto de saída para testar uma máquina em produção?
  \item Por que a edição \emph{on-line} tem efeito imediato, e que cuidado isso exige?
  \item Um painel apresenta falhas que só acontecem nos meses mais quentes. Que modo de falha
        suspeitar e como confirmar?
  \item Por que é obrigatório testar a restauração de um backup, e não apenas produzi-lo?
  \item Escolha três itens da Tabela~\ref{tab:rotina-preventiva} e diga qual evidência prova que
        foram executados.
  \item Cite duas fontes de dado já disponíveis que permitem manutenção preditiva sem sensor novo, e
        diga o que cada uma antecipa.
  \item Em que condições a manutenção preditiva não compensa? Dê duas condições.
  \item Um fornecedor propõe um pacote de predição que ``avisa tudo automaticamente''. Que perguntas
        você faz antes de assinar? Relacione com os três requisitos da predição viável.
\end{enumerate}
\end{exercicios}
